% Optional main-text insertion. The defining result and complete proof are in
% sections/posterior_certificates_theory.tex; do not duplicate theorem numbering.
\paragraph{Posterior decisions and logical certificates.}
A tournament law and a rule for reporting its UC are separate objects. GFM uses the joint moments $\Pr(i\in\UC(T),|\UC(T)|=s\mid D)$ to maximize modeled expected F1 \citep{waegeman2014fmeasure}; membership marginals alone omit its size-dependent denominator. Shared mixture components also induce dependent edges, so multiplying mixture edge marginals generally gives a different law.

Known orientations can alternatively bound UC without recovering every edge. Define $I(G)$ as vertices with known paths of length at most two to every opponent, and exclude $v$ from $O(G)$ when a known $u\to v$ has, for each $w\notin\{u,v\}$, known $w\to v$ or $u\to w$. Proposition~\ref{prop:uc-completion-enclosure} proves $I(G)\subseteq\UC(T)\subseteq O(G)$ for every strict completion $T$; equality certifies the complete set. This applies established two-step tournament supports \citep{aziz2015partial,contet2026supports}. Fixed-design i.i.d. pair intervals make the enclosure simultaneous with probability $1-\delta$ (Corollary~\ref{cor:uc-fixed-design-enclosure}). That sampling claim does not apply to finite human-cohort subsampling, and the certificate establishes no adaptive-query or comparison-budget savings. Appendix~\ref{app:posterior-certificates-theory} gives complete proofs and the separate mixture-training identity.
